% Copyright© Maia Japhoshvili
\documentclass[11pt]{book}

\textwidth=165mm
\oddsidemargin=-4.4mm
\evensidemargin=4.4mm
\textheight=235mm
\topmargin=-5mm
\headheight=5.0mm
\headsep=6.6mm
\footskip=8.1mm

\usepackage{times,fancyhdr}
\usepackage{amssymb,amsbsy,exscale,times}
\usepackage{amsmath,amssymb,amstext,amscd,amsfonts}
\usepackage{theorem,graphicx}
\usepackage{mathrsfs}
\usepackage{color}
\usepackage{index}
\usepackage{verbatim}
\usepackage{url}

\makeindex
\raggedbottom


\begin{document}


\addcontentsline{toc}{section}{\textbf{Besik Dundua, Furio Honsell, Mariam Katamashvili, David Maisuradze}, Fuzzifying Derivatives of Regular Expressions \medskip}
\index{Dundua Besik}
\index{Honsell Furio}
\index{Katamashvili Mariam}
\index{Maisuradze David}


\begin{center}
    \textsf{\textbf{\Large Fuzzifying Derivatives of Regular Expressions}} \\[0.5cm]
    \textbf{Besik Dundua$^{1,2}$, Furio Honsell$^{3}$, Mariam Katamashvili$^{1,2}$, David Maisuradze$^{1}$} \\[0.2cm]
    $^{1}$\textit{Kutaisi International University, Kutaisi, Georgia} \\[0.05cm]
    $^{2}$\textit{VIAM, Tbilisi State University, Tbilisi, Georgia} \\[0.05cm]
    $^{3}$\textit{University of Udine, Udine, Italy}
\end{center}

\bigskip

There is growing interest in computer science and artificial intelligence in moving
from {\em exact}, {\em qualitative} settings to {\em quantitative}, and ultimately
{\em fuzzy}, ones. In this
spirit, we extend the theory of finite automata, based on exact symbol equality, to
fuzzy symbol matching. Following
similarity-based reasoning~[7], we enrich the alphabet with a {\em similarity relation}
$\mathcal R$ and perform matching up to a fixed {\em cut} value $\mu\in(0,1]$: two
symbols match when their similarity is at least $\mu$. Here $\mathcal R$ takes values in
$[0,1]$ endowed with a left-continuous T-norm (a residuated lattice~[4]), and is
reflexive, symmetric, and transitive.

We generalize the theory of Brzozowski derivatives~[1], revisited in~[5], which
constructs automata directly from regular expressions: after reading a word, the
current state is the corresponding derivative of the expression. The similarity relation
on symbols is extended to regular expressions, words, and languages, yielding in each
case a similarity relation. Since intersection and complement do not fuzzify well, our
results concern the {\em restricted regular expressions}: empty language, empty word,
letters, union, concatenation, and Kleene star.

We introduce {\em fuzzy language and expression derivatives}, in which each symbol of the
input word is matched up to the cut value $\mu$, and prove that they preserve
regularity. We further establish a
{\em semantic-syntactic correspondence} (the language of the fuzzy derivative of an
expression coincides with the fuzzy derivative of its language) and a {\em finiteness}
result: iterating fuzzy derivatives over all words yields only finitely many distinct
languages.

One of the main results connects our fuzzy theory to the classical, crisp one. For a
fixed cut value $\mu$ and G\"odel's minimum T-norm, the similarity relation induces an
equivalence on the alphabet, whose quotient we call the {\em quotient alphabet}
$\Sigma^\mu$. We prove that ordinary
Brzozowski derivative automata over $\Sigma^\mu$ capture exactly the fuzzy derivatives
over the original alphabet. Most of these results are formalized in the Rocq prover~[6], see~[3],
building on the Coquand-Siles~[2] formalization of regular-expression equivalence via
Brzozowski derivatives. Furthermore, we explore how to relate our similarity-based
approach to the fuzzy-set construction of Stamenkovi\'c and \'Ciri\'c~[8].

\subsection*{References}

\begin{enumerate}

\item[{[1]}] J. A. Brzozowski, Derivatives of regular expressions. \emph{J. ACM} \textbf{11} (1964), no. 4, 481--494.

\vspace{-7pt}
\item[{[2]}] T. Coquand and V. Siles, A decision procedure for regular expression equivalence in type theory. \emph{Proc. 1st Int. Conf. on Certified Programs and Proofs (CPP)}, 2011.

\vspace{-7pt}
\item[{[3]}] M. Katamashvili and D. Maisuradze, Rocq formalization for regular-expression derivatives with similarity. \url{https://github.com/DUT1K1/regular-expression-derivative}.

\vspace{-7pt}
\item[{[4]}] G. Metcalfe, Fundamentals of fuzzy logics. Lecture notes, Vienna University of Technology, 2005.

\vspace{-7pt}
\item[{[5]}] S. Owens, J. Reppy and A. Turon, Regular-expression derivatives re-examined. \emph{J. Funct. Programming} \textbf{19} (2009), no. 2, 173--190.

\vspace{-7pt}
\item[{[6]}] The Rocq Prover Development Team, The Rocq Prover, version 9.1.0. \url{https://rocq-prover.org/releases/9.1.0}, 2025.

\vspace{-7pt}
\item[{[7]}] M. I. Sessa, Approximate reasoning by similarity-based SLD resolution. \emph{Theoret. Comput. Sci.} \textbf{275} (2002), no. 1--2, 389--426.

\vspace{-7pt}
\item[{[8]}] A. Stamenkovi\'c and M. \'Ciri\'c, Construction of fuzzy automata from fuzzy regular expressions. \emph{Fuzzy Sets and Systems} \textbf{199} (2012), 1--27.
\end{enumerate}


\end{document}
